% ============================================================
% Project: THE PHILOSOPHY OF VIABILITY
% Book: Book I — Ontology of Viability
% Book ID: VS.B01
% Part: Part I — Identity and Manifestation
% Chapter: Chapter 07 — Formal and Logical Reconstruction
% File: 2_MAINMATTER/PART_I_IDENTITY_AND_MANIFESTATION/CHAPTER_07_FORMAL_AND_LOGICAL_RECONSTRUCTION/sections/section_03_formal_domains_and_types.tex
% Section: 03 — Formal Domains and Types
% Purpose: To establish the distinctions among domains of
% discourse, formal domains of interpretation, and ontological
% categories; to specify the roles and type constraints of
% formal expressions without converting formal architecture
% into unsupported ontological commitments.
% ============================================================

\section{Formal Domains and Types}
\label{sec:ch07_s03}

\subsection{The Problem of Domains in Ontological Formalization}

The formal reconstruction of an ontological framework requires
an account of the domains over which its expressions range.
Variables must range over something; predicates must apply
to specified kinds of terms; relations must connect arguments
of identifiable types; and functions, where introduced,
must have determinate input and output conditions.

These requirements are formal, but their consequences are
philosophically significant. The selection of a domain may
silently determine what the formal system is capable of
representing. The selection of types may suggest divisions
among the entities under consideration. The use of a
particular predicate may introduce distinctions that were
not established in the ontology itself.

The central methodological problem is therefore how to
specify formal domains and types while preserving the
distinction between the organization of a representational
system and the ontological organization of its subject matter.

A formal domain is not automatically an ontological realm.
A type is not automatically an independently existing
kind of being. A variable range does not, merely by being
specified, establish the complete extension of an
ontological category.

At the same time, formal domains cannot be selected
arbitrarily. Their adequacy depends upon the subject
matter and claims that the formal reconstruction is
intended to represent. A formal language that cannot
distinguish existents from expressions, or relations
from properties, may be inadequate for particular
arguments even if it is technically well-formed.

The task is consequently twofold: to establish a
formally controlled architecture and to make explicit
the interpretive commitments through which that
architecture is connected to ontology.

\subsection{Three Senses of Domain}

The term \emph{domain} is frequently used in formal
and philosophical discussions without sufficient
qualification. At least three senses must be
distinguished in the present reconstruction.

\paragraph{Domain of discourse.}

The domain of discourse specifies what a particular
discussion is about. It establishes the intended
subject matter and the range of items relevant to
the argument. A discussion may concern existents
in general, a particular class of existents, relations,
manifestations, or formal expressions.

A domain of discourse is therefore determined by
the question under examination. Its specification
does not necessarily provide a complete formal
semantics, and it does not establish that the
items within it form an ontologically independent
category.

\paragraph{Formal domain of interpretation.}

A formal domain of interpretation is the collection
of objects over which the variables of a formal
language range under a specified interpretation.
It is part of the semantic structure assigned to
the language.

In a standard first-order interpretation, the
domain is ordinarily taken to be non-empty. This
is a condition of that familiar semantic framework.
It must not be mistaken for an ontological
principle that reality necessarily contains an
object because a particular formal semantics
requires a non-empty domain.

Nor does the use of one formal domain establish
that every expression in the language must refer
to an independently existing object. Predicates,
relations, and functions receive their interpretations
according to the semantic rules of the system.

\paragraph{Ontological domain.}

An ontological domain is a proposed range of
what exists or of a specified ontological category.
A claim that a particular kind of entity constitutes
an ontological domain is therefore a substantive
claim about reality, not merely a decision about
the interpretation of a formal language.

The ontological legitimacy of such a domain
requires philosophical justification. It cannot
be inferred solely from the fact that a formal
system introduces a corresponding type or
variable range.

These three senses may be related in a particular
reconstruction, but they must not be conflated.
A domain of discourse determines the subject
of an argument; a formal domain determines
the range of formal variables under an
interpretation; an ontological domain asserts
something about the organization of what exists.

The distinction is represented in
Figure~\ref{fig:ch07_domain_architecture}.

% =============================================================================
% Project : THE_PHILOSOPHY_OF_VIABILITY
% Book    : Book I — Ontology of Viability
% Book ID : VS.B01
% Part    : Part I — Identity and Manifestation (VS.B01.P01)
% Chapter : Chapter 07 — Formal and Logical Reconstruction
% Chapter ID : VS.B01.P01.C07
%
% File    : 4_ASSETS/diagrams/VS.B01.P01.C07.S03.SS02.P001.DG01.domain_architecture_and_interpretation.tex
%
% Section / Chapter / Asset:
%   Diagram DG01 — Domain Architecture and Interpretation
%
% Structural Identifier : VS.B01.P01.C07.S03.SS02.P001.DG01
% Asset Type : DG
%
% Purpose :
%   To distinguish the domain of discourse, the formal domain of
%   interpretation, and the ontological domain without identifying their
%   status.
%
% Representational limit :
%   The arrows represent methodological and interpretive relations, not
%   ontological production or equivalence.
%
% Dependencies :
%   TikZ; libraries used: arrows.meta, positioning. All are loaded by the
%   project preamble; no per-file package or library loading.
%
% =============================================================================

\begin{figure}[htbp]
    \centering
    \hyphenpenalty=10000 \exhyphenpenalty=10000\relax  % no hyphenation inside diagram labels

    \begin{tikzpicture}[
        font=\small,
        node distance=1.4cm,
        % --- base box ---
        box/.style={
            draw, rounded corners=3pt, align=center,
            text width=8.2cm, minimum height=1.9cm,
            inner sep=7pt, line width=0.8pt
        },
        % --- node roles (names avoid clashes with built-in TikZ keys) ---
        discbox/.style={box, fill=blue!8,  draw=blue!55!black},
        formbox/.style={box, fill=blue!16, draw=blue!55!black},
        ontbox/.style={box, dashed, fill=gray!10, draw=black!60},
        % --- arrows (methodological and interpretive, not production) ---
        directed/.style={-{Stealth[length=2.5mm]}, line width=0.9pt},
        qualified/.style={{Stealth[length=2.5mm]}-{Stealth[length=2.5mm]},
                          line width=0.9pt, dashed},
        elabel/.style={font=\footnotesize\itshape, text=black!75,
                       align=left, text width=3.4cm}
    ]

    % ---------------------------------------------------------------
    % Three senses of domain (stacked; no identity of status asserted)
    % ---------------------------------------------------------------
    \node[discbox] (discourse) {
        \textbf{Domain of Discourse}\\[3pt]
        Subject matter selected for a particular inquiry\\
        Existent, identity, determination, relation, manifestation
    };
    \node[formbox, below=of discourse] (formal) {
        \textbf{Formal Domain of Interpretation}\\[3pt]
        Variable ranges, formal objects, assigned interpretations\\
        Domains, sorts, predicates, relations, functions
    };
    \node[ontbox, below=of formal] (ontological) {
        \textbf{Ontological Domain}\\[3pt]
        A proposed range or category of what exists\\
        Requires independent ontological justification
    };

    \draw[directed] (discourse) -- (formal)
        node[elabel, midway, right=5pt] {guides formal representation};
    \draw[qualified] (formal) -- (ontological)
        node[elabel, midway, right=5pt] {interpretation may refer to ontological claims};

    % ---------------------------------------------------------------
    % Legend and note
    % ---------------------------------------------------------------
    \node[below=0.6cm of ontological, font=\scriptsize] (legend) {
        \tikz[baseline=-0.5ex]\draw[directed] (0,0) -- (0.7,0);~methodological guidance
        \quad\quad
        \tikz[baseline=-0.5ex]\draw[qualified] (0,0) -- (0.7,0);~qualified interpretive relation
    };

    \node[below=0.4cm of legend, text width=10cm, align=center,
          font=\footnotesize\itshape] {
        Formal domains and sorts do not, by themselves,
        establish the divisions or completeness of reality.
    };

    \end{tikzpicture}

    \caption{Three senses of domain in formal ontological reconstruction.
    The arrows represent methodological and interpretive
    relations, not ontological production or equivalence.}
    \label{fig:ch07_domain_architecture}
\end{figure}

% =============================================================================
% END OF DIAGRAM DG01
% =============================================================================


The diagram presents a methodological relation
among the three senses. The domain of discourse
guides the selection of a formal representation,
while an ontological domain, where asserted,
requires independent justification. The diagram
does not imply that the formal domain creates
or exhausts the ontological domain.

\subsection{Formal Domains and the Ontological Subject Matter}

The principal subject matter of Part I is the
ontological examination of identity and its
relations to being, determination, distinction,
unity, plurality, invariance, expression,
manifestation, persistence, and continuity.

The formal reconstruction must be capable of
representing claims concerning these matters.
It does not follow, however, that every concept
must be represented as an object in one
undifferentiated domain.

For example, the existent under consideration
may be represented by an object variable.
A property attributed to that existent may
be represented by a predicate. A relation
between two existents may be represented
by a relational expression. A transformation
may be represented through a relation between
conditions or states, or through a function
where the relevant conditions of functional
representation have been justified.

These representational choices do not establish
that properties, relations, transformations,
and states are ontologically independent
objects. They establish only that the formal
language uses different expression types
to articulate different logical roles.

The distinction is particularly important
when formalizing identity. If identity is
represented by a predicate, the reconstruction
must clarify what the predicate expresses.
If numerical identity is represented by the
logical equality symbol, that symbol must
retain its established logical function.
If identity as an ontological concern is
represented through a more elaborate
formal structure, the interpretation of
that structure must be explicitly justified.

A single formal domain may be adequate for
some arguments. Other arguments may require
a more differentiated language. The choice
between these possibilities is a matter
of formal adequacy and interpretive purpose,
not an immediate discovery of the number
or structure of ontological domains.

\subsection{Types as Formal Constraints}

A \emph{type} specifies the formal role or
category of an expression within a language.
It restricts how that expression may be used
and which combinations of expressions are
well-formed.

Types can prevent errors that arise when
expressions with different roles are treated
as interchangeable. A variable intended to
range over existents, for example, should
not be used without qualification as though
it necessarily ranged over propositions,
formal symbols, or relations.

A type system may also distinguish different
kinds of arguments. A binary relation may
require two arguments of the same type,
or it may connect arguments of different
types. A function may accept an argument
of one type and return a value of another.
Such restrictions make the formal language
more precise.

However, a formal type is not automatically
an ontological classification. The type
assigned to an expression describes its
role within the formal system. Whether
that role corresponds to a real ontological
distinction is a further question.

The reconstruction must therefore distinguish
two claims:

\begin{enumerate}
    \item \textbf{Formal typing claim:}
    an expression belongs to a specified
    formal type and may be used only in
    accordance with the rules governing
    that type.

    \item \textbf{Ontological classification claim:}
    the subject matter represented by
    that expression belongs to a specified
    ontological category.
\end{enumerate}

The first may be established through the
definition of the formal language. The
second requires ontological justification.

This distinction allows the formal system
to use types as instruments of precision
without turning every notational distinction
into a metaphysical division.

\subsection{A Provisional Typing Architecture}

For the purposes of the present reconstruction,
a provisional distinction among formal expression
types is useful. The following categories
identify logical roles rather than independent
ontological substances.

\paragraph{Object variables and constants.}

Object variables range over the objects of
a specified formal domain. Constants designate
particular elements of that domain under
an interpretation.

When the domain is intended to represent
existents, object variables may be used to
refer to existents under examination.
The notation alone does not determine
whether the domain includes every existent
or only those relevant to a particular
argument.

\paragraph{Predicates.}

A predicate expresses a condition or
characteristic attributed to one or more
arguments. A unary predicate applies to
one argument; a predicate of higher arity
applies to multiple arguments.

For example, a predicate may be used
to represent a specified condition of
an existent. The use of a predicate
does not by itself establish whether
the condition is essential, accidental,
identity-constitutive, or merely descriptive.
Those interpretations require further
ontological specification.

\paragraph{Relations.}

A relation connects arguments under
a specified formal interpretation.
Its arity and argument types must be
identified. A relation may connect
two existents, an existent and a
condition, or other formally specified
items.

The use of a relation symbol does not
settle whether the relation is internal,
external, constitutive, derivative,
or independent. These are interpretive
and ontological questions that cannot
be resolved by arity alone.

\paragraph{Functions.}

A function assigns an output to an
input or tuple of inputs according
to a specified rule. In a formal
system, the introduction of a function
normally carries a requirement that
the relevant input be associated
with a unique output in the function's
intended domain.

This requirement can be too strong
for some ontological relations.
Where a process or transformation
does not yield a unique result under
the intended conditions, a relation
or a suitably restricted partial
function may be more appropriate.

The formal choice must reflect the
claim being represented. It must not
impose uniqueness merely because
a functional notation is convenient.

\paragraph{Operators.}

An operator acts upon expressions
of a specified formal category and
produces another expression or value.
Its use requires a defined input
and output structure. An operator
introduced for formal manipulation
must not be mistaken for an ontological
agent or process.

\paragraph{Transformations and transition relations.}

A transformation may be represented
as a function when each relevant
input determines a unique output.
It may instead be represented as
a relation when multiple outcomes,
partial applicability, or alternative
successors must be allowed.

A transition relation can represent
a connection between two conditions
without asserting that the connection
is deterministic or reversible.

These formal possibilities are
important for later examination
of invariance, change, persistence,
and continuity. They do not
predetermine the ontological
interpretation of those phenomena.

\paragraph{Claims and propositions.}

A proposition is a content-bearing
expression that may be asserted,
denied, supported, or left unresolved.
Its logical role differs from that
of an object variable or predicate.
The formal language must therefore
distinguish the expressions used
to refer to objects from those
used to formulate claims about them.

The principal formal expression
types and their constraints are
summarized in Table~\ref{tab:ch07_formal_types}.

% ============================================================
% Project: THE PHILOSOPHY OF VIABILITY
% Book: Book I — Ontology of Viability
% Book ID: VS.B01
% Chapter: Chapter 07 — Formal and Logical Reconstruction
% File: 4_ASSETS/tables/VS.B01.P01.C07.S03.SS04.P001.TL01.formal_domains_and_types.tex
% Asset: TL01 — Formal Domains and Types
% Purpose: Summary of formal expression types and their
% interpretive and ontological constraints.
% ============================================================

\begin{table}[htbp]
\centering
\small
\caption{Formal Expression Types and Their Roles}
\label{tab:ch07_formal_types}
\renewcommand{\arraystretch}{1.25}
\begin{tabularx}{\textwidth}{|p{0.19\textwidth}|p{0.27\textwidth}|X|}
\hline
\textbf{Formal Type} &
\textbf{Formal Role} &
\textbf{Interpretive and Ontological Constraint}
\\
\hline
Object variable &
Ranges over a specified formal domain. &
Its range must be stated; it does not establish the complete inventory of existents.
\\
\hline
Constant &
Designates an element under an interpretation. &
Designation within a formal model does not independently establish ontological status.
\\
\hline
Predicate &
Expresses a condition attributed to one or more arguments. &
Formal predication does not establish whether the attributed condition is essential or identity-constitutive.
\\
\hline
Relation &
Connects arguments of specified types and arity. &
Its formal properties must not be transferred automatically to the represented ontological relation.
\\
\hline
Function &
Assigns a unique output to each admissible input. &
Totality and uniqueness must be justified; otherwise a partial function or relation may be required.
\\
\hline
Operator &
Acts upon formal expressions according to specified rules. &
Formal operation must not be mistaken for an ontological agent or process.
\\
\hline
Transformation &
Represents a transition between conditions or configurations. &
Must specify its inputs, outputs, domain, and whether it is functional or relational.
\\
\hline
Proposition &
Expresses content that may be asserted, denied, or assessed. &
Must be distinguished from the objects and expressions about which it makes claims.
\\
\hline
\end{tabularx}
\end{table}

% ============================================================
% END OF ASSET TL01 — CHAPTER 07, SECTION 03
% ============================================================

The table is not a final ontological
taxonomy. It is a working specification
of formal roles that will be refined
when the subsequent sections define
the vocabulary, commitments, and
inference rules of the reconstruction.

\subsection{Single-Sorted and Many-Sorted Formal Languages}

A single-sorted formal language uses
one general domain for its object
variables. Its expressions may
nevertheless distinguish properties,
relations, and other roles through
predicates and relational symbols.

Such a language can be economical
and may be adequate for arguments
that do not require explicit typing
of different object categories.
Its simplicity, however, can make
some distinctions less visible and
can permit combinations of expressions
that a more restricted language
would exclude.

A many-sorted language distinguishes
two or more formal sorts. Variables
are assigned to specified sorts,
and predicates, relations, and
functions are given corresponding
sort signatures.

For example, a language may
distinguish variables intended
to range over existents from
variables intended to range
over formal expressions. This
distinction can prevent a formal
expression from being treated
as an existent merely because
both are represented by variables
in a general language.

Many-sorted formalization is
therefore a useful candidate
for parts of the present
reconstruction. It permits
explicit control over the
argument types of relations
and the input and output types
of functions.

Nevertheless, a many-sorted
language introduces a further
methodological responsibility.
The sorts must be justified
as formal requirements of
the intended representation.
They must not be presented
as proof that reality consists
of a corresponding number
of independent ontological
categories.

A sort may be introduced
for technical convenience,
for semantic clarity, or
because a particular argument
requires a distinction.
These grounds are not
equivalent to the claim
that the sort represents
a fundamental division
in being.

The choice between single-sorted
and many-sorted languages
must therefore be made
locally, according to the
requirements of each argument.
No universal preference
is established at this stage.

\subsection{The Domain of Existents and the Question of Quantification}

Since Part I concerns identity
as an ontological issue, the
formal reconstruction will
frequently require variables
whose intended range consists
of existents.

For convenience, let \(D_E\)
denote a formal domain intended
to represent the existents
relevant to a particular
interpretation. A variable
\(x\) may then be assigned
the range:

\[
x \in D_E
\]

This notation expresses
a formal typing or range
restriction. It does not
establish that \(D_E\)
is a complete inventory
of everything that exists.
Nor does it imply that
all existents are exhaustively
captured by one formal
representation.

The domain may be restricted
to the subject matter of
a particular argument.
A formal reconstruction
of a specific proposition
need not quantify over
every possible existent.
The scope of quantification
must therefore be made
explicit whenever it affects
the interpretation or
validity of a claim.

It is also necessary
to distinguish quantification
over existents from
quantification over expressions
or claims. If a language
contains variables for
formal expressions, those
variables belong to a
different formal role.
The existence of a formal
expression within a language
does not make it an existent
in the ontological sense
under examination.

This distinction becomes
especially important when
the reconstruction refers
to definitions, propositions,
models, or proofs. These
are components of the
formal inquiry and must
not be silently included
in the same ontological
range as the existents
about which the inquiry
makes claims.

\subsection{Predicates, Properties, and Ontological Attribution}

Predicates are among
the principal means
through which formal
languages represent
attribution. A unary
predicate may express
that an existent satisfies
a specified condition.
A relational predicate
may express that two
or more arguments stand
in a specified relation.

The formal distinction
between a predicate
and a property must
nevertheless be preserved.

A predicate is an
expression of a formal
language. A property,
when used ontologically,
concerns what is attributed
to or belongs to an existent.
A predicate may be
interpreted as representing
a property, but the
formal expression and
the ontological feature
are not thereby identical.

This distinction prevents
a semantic decision
from becoming an
unexamined metaphysical
commitment. A formal
system may interpret
a predicate extensionally,
through the objects
to which it applies,
without thereby settling
the ontological status
of properties.

For example, suppose
\(P(x)\) is used to
express that an existent
satisfies a specified
condition. The formal
expression permits
the formulation:

\[
P(x)
\]

It does not, without
additional interpretation,
establish that the
condition is essential
to \(x\), that it
constitutes the identity
of \(x\), or that it
is an independently
existing property.

The ontological force
of the attribution
must be established
by the argument in
which it occurs.

The same restriction
applies to predicates
used for determination,
unity, manifestation,
persistence, or continuity.
The presence of a
predicate in a formal
language does not
settle the precise
ontological status
of the phenomenon
it represents.

\subsection{Relations and Their Argument Types}

The formal representation
of relations requires
more than the introduction
of a relation symbol.
Its arguments, arity,
and intended interpretation
must be specified.

A binary relation
may be represented
schematically as:

\[
R(x,y)
\]

The expression is
formally meaningful
only when the language
specifies the types
of \(x\) and \(y\),
the interpretation
of \(R\), and the
conditions under
which the relation
is satisfied.

A relation between
two existents is
not necessarily of
the same kind as
a relation between
an existent and
a condition, or
between two formal
expressions. If the
formal language does
not distinguish these
roles where the
argument requires
it, ambiguities
may arise.

The arity of a
relation does not
establish its ontological
character. A binary
relation may be
symmetric, asymmetric,
reflexive, transitive,
or possess none
of these properties.
Such features must
be established
for the relation
under examination.

Nor should the
formal properties
of a relation
be transferred
automatically to
the ontological
relation it is
intended to represent.
The adequacy of
that interpretation
requires separate
analysis.

For the present
reconstruction, every
relation introduced
to represent an
ontological claim
must therefore be
accompanied by
an explicit account
of its argument
types and intended
meaning.

\subsection{Functions, Partiality, and Non-Deterministic Relations}

Functions provide
a compact means
of representing
determinate mappings.
However, they carry
formal conditions
that may not be
appropriate for
every ontological
process or relation.

A function from
a domain \(A\)
to a domain \(B\)
assigns a unique
output in \(B\)
to each admissible
input in \(A\).
Its formal signature
may be written:

\[
f : A \rightarrow B
\]

The notation implies
that the function
is defined for
every element of
its stated domain
and that each
input has one
unique output.

If an ontological
process can produce
more than one
possible outcome,
or if it is not
defined for every
relevant input,
a total function
may impose an
unjustified restriction.

Three possibilities
must then be
considered.

First, the
function's domain
may be restricted
to the inputs
for which the
mapping is defined.

Second, a partial
function may be
used where the
formal framework
supports partiality.

Third, a relation
may be used
when one input
can be connected
to several possible
outputs.

These alternatives
are formally distinct.
Their selection
must reflect the
structure of the
claim being represented.

This distinction
will become
important when
the reconstruction
examines transformation
and invariance.
A transformation
must not be
assumed to be
deterministic merely
because it is
represented through
a function. Conversely,
the use of a
relation does not
by itself establish
ontological indeterminacy.
It may simply
reflect the limits
or intended scope
of the formal
representation.

\subsection{Type Signatures and Formal Well-Formedness}

A type signature
specifies the formal
types of the
arguments accepted
by an expression
and, where relevant,
the type of
its output.

For a predicate
or relation, a
signature may
indicate the
types of its
arguments. For
a function, it
indicates the
input and output
types.

A signature
therefore establishes
conditions of
formal well-formedness.
It helps determine
whether an expression
is admissible
within the language.

The signature
of a binary
relation may
be represented
schematically as:

\[
R : A \times B
\]

This notation
indicates that
the first argument
belongs to type
\(A\) and the
second to type
\(B\). It does
not, by itself,
assert that \(A\)
and \(B\) are
ontologically
independent domains.

Similarly, a
function signature:

\[
f : A \rightarrow B
\]

specifies a
formal mapping
between types.
It does not
establish that
the mapping
corresponds to
a causal process,
an ontological
production, or
a necessary
relation in
reality.

Type signatures
must consequently
be treated as
formal constraints
whose ontological
interpretation
remains subject
to argument.

\subsection{The Formal Identity Relation and Type Restrictions}

The logical identity
relation requires
particular attention
because it is
often used as
a basic formal
resource.

Within a single-sorted
first-order language,
the expression
\(x=y\) ordinarily
states that \(x\)
and \(y\) designate
the same object
under the interpretation.

In a many-sorted
language, the
use of identity
must respect
the formal
framework's
typing conventions.
Identity is
normally meaningful
between terms
that designate
objects of the
same sort, unless
a more elaborate
system explicitly
provides another
interpretation.

This restriction
is a matter
of formal
well-formedness.
It must not
be confused
with a claim
that ontological
identity is
constituted by
formal sort
membership.

The formal
identity relation
continues to
express numerical
identity within
the interpretation.
It does not
supply a complete
account of
ontological identity
as developed
in Part I.

Consequently,
the formal
reconstruction
must maintain
the distinction
between the
logical relation
of being one
and the same
object and
the ontological
question concerning
what it means
for an existent
to be the
existent that
it is.

\subsection{Formal Domains Must Not Predetermine Ontological Plurality}

The introduction
of several formal
domains or sorts
may be useful
for representing
different kinds
of arguments.
It does not
establish that
the ontology
contains the
same number
of independent
kinds of being.

Likewise, the
use of one
formal domain
does not
establish that
all existents
are ontologically
homogeneous.

Formal plurality
and ontological
plurality must
be examined
separately.

A language
may distinguish
objects, relations,
states, and
expressions because
these play
different roles
in formal
representation.
It does
not follow
that each
role corresponds
to an independent
ontological
entity or
category.

Conversely,
a single
formal sort
may contain
items that
are ontologically
different in
ways not
represented by
the type system.

The number
of formal
sorts is
therefore not
a measure
of the number
of identities,
the number
of ontological
categories,
or the
number of
independent
realities.

This restriction
is particularly
important for
the treatment
of unity
and plurality.
The formal
classification
of components,
relations,
manifestations,
or conditions
must not
be converted
into a conclusion
about the
plurality of
identities without
a separate
argument.

\subsection{The Choice of Formal Architecture}

No single
formal architecture
is established
as universally
adequate for
every claim
in the present
ontology.

A single-sorted
first-order language
may be sufficient
for certain
arguments involving
existents and
relations. A
many-sorted language
may be more
suitable where
the distinction
between existents,
formal expressions,
and other
argument types
is necessary.

A relational
representation
may be preferable
to a functional
one where
multiple outcomes
or partial
applicability must
be preserved.
A richer
formal language
may be required
where the
argument involves
modal, temporal,
or other
structures that
cannot be
adequately represented
by the selected
first-order resources.

These possibilities
must be evaluated
according to
the claims
to be represented,
the distinctions
that must be
preserved, and
the inferential
consequences of
the formal
choices.

The choice
of a formal
system is
therefore provisional
until its
adequacy has
been examined
against the
relevant ontological
arguments.

The reconstruction
must not select
a formal
architecture merely
because it
is familiar,
compact, or
mathematically
convenient. Nor
should it
introduce complexity
that does
not serve
a specific
analytical purpose.

\subsection{Open Questions Concerning Domains and Types}

Several questions
remain open
at this stage.

It remains
to be determined
whether the
formal reconstruction
of identity
can be carried
out adequately
within one
general domain
or whether
particular arguments
require distinct
sorts.

It remains
to be determined
how determination
should be
represented in
different contexts:
as a predicate,
a relation,
a structured
configuration,
or a combination
of these.

It remains
to be determined
whether the
formal representation
of manifestation
and expression
requires distinct
types or
whether their
difference can
be adequately
represented through
distinct predicates
and relations
within a
shared domain.

It remains
to be determined
which transformations
should be
represented as
functions and
which require
relational or
partial representations.

It also remains
to be determined
whether the
formal distinctions
required for
the ontology
can be
represented without
introducing unnecessary
ontological commitments
through the
type system itself.

These questions
must be addressed
in the subsequent
formal development.
Their present
open status
is not a
reason to
avoid formalization;
it is a
reason to
make the
assumptions of
formalization explicit.

\subsection{Conclusion: Domains and Types as Instruments of Precision}

This section
has distinguished
the domain
of discourse,
the formal
domain of
interpretation, and
the ontological
domain. It
has established
that these
senses of
domain are
related but
not interchangeable.

It has
also distinguished
formal types
from ontological
categories and
identified the
principal expression
roles required
for the
reconstruction:
object variables,
constants,
predicates,
relations,
functions,
operators,
transformations,
and propositions.

The analysis
has shown
that the
choice of
formal domains
and types
must be
guided by
the ontological
claims under
examination, while
remaining distinct
from the
justification of
those claims.

A formal
architecture is
adequate only
to the extent
that it
preserves the
relevant distinctions,
makes its
interpretive commitments
explicit, and
does not
derive ontological
conclusions merely
from its
own technical
organization.

The next
section examines
definitions and
conceptual dependencies.
It will determine
how the terms
introduced in
the preceding
section can
be defined,
which concepts
depend upon
others for
their articulation,
and how
circularity or
unjustified
conceptual reduction
can be
identified within
the reconstructed
framework.

% ============================================================
% END OF SECTION 03 OF CHAPTER 07
% ============================================================